\section{Registration of the GSLC pairs}
\label{app:registration}
The PR GSLCs of the two dates are each geocoded with their own ionospheric (TEC) and solid-earth-tide geolocation
corrections. Those corrections differ between dates by about a resolution cell (2.5--3.2\,m on the Panama D141 pair),
so an interferogram formed directly from two GSLCs loses about half its coherence. We keep the first date's GSLC and
geocode the second date's RSLC onto its grids. That geocoding drops the second date's own corrections and adds a
registration field: the first date's correction shift plus the offsets measured between the two dates (our port of
ISCE3's ROFF correlator). The side band is geocoded with the same field in its own pixels; the TEC part of the shift
scales by $(\fA/\fB)^2$. Both bands of both dates then lie on common grids, and the cells of
Section~\ref{sec:estimator} are sums over identical pixels.

\section{ISCE3's correction passes, as ported}
\label{app:isce3}
With $\mathcal{G}$ a NaN-aware truncated Gaussian ($\sigma$ 2.0\,km by 2.6\,km, truncated at 1.75 and 1.49 $\sigma$:
ISCE3's 14 by 110 taps for $\sigma$ 4 by 37 side-band pixels), the first pass
(\texttt{unwrapping\_correction\_with\_filter}, one iteration) replaces every unwrapped field $u$ ($\phiA$ and
$\Delta$, or $\phiA$ and $\phiB$) by
\begin{equation}
  u \leftarrow u - 2\pi \left\lfloor \frac{u - \mathcal{G}u}{2\pi} \right\rceil .
\end{equation}
The second pass (\texttt{compute\_unwrapp\_error\_main\_diff\_ms\_band}) forms the first estimates
$\hat d = c_0\phiA + c_1\Delta$ and $\hat n = \phiA - \hat d$, filters both with the dispersive filter ($F$), and
corrects
\begin{equation}
  \phiA \leftarrow \phiA - 2\pi\left\lfloor\frac{\phiA - F\hat n - F\hat d}{2\pi}\right\rceil,\qquad
  \Delta \leftarrow \Delta - 2\pi\left\lfloor\frac{\Delta + \frac{\Delta f}{\fA}F\hat n - \frac{\Delta f}{\fB}F\hat d}
  {2\pi}\right\rceil ,
\end{equation}
before the final estimate and filter. The two-band version
(\texttt{compute\_unwrapp\_error\_main\_side\_band}) rounds the differential integer
$\lfloor[(1-\fB/\fA)F\hat n + (1-\fA/\fB)F\hat d + \phiB - \phiA]/2\pi\rceil$ and the common one from the sum, as in
\citet[eqs.~28--29]{gomba2016toward}. ISCE2's ALOS-2 processor instead rounds $\phiA-\phiB$ against a weighted
quadratic of itself \citep{liang2018insar}. These are the \texttt{diff\_isce3}, \texttt{side\_isce3} and
\texttt{side\_isce2} estimators of the experiments.

\section{The synthetic generator}
\label{app:synthetic}
The scene is 96 by 96\,km of 40\,m pixels. The truth is a vertical left-lateral rupture striking N105$^\circ$E,
built from Okada rectangles \citep{okada1985surface}: 10 segments along strike by two in depth (0--10 and 10--15\,km;
the lower layer carries 0.6 of the slip), with slip tapered along strike as $\sqrt{\sin}$. It is 50\,km long and ends
inside the scene (``finite''), or 160\,km long and runs through it (``through''). The LOS is a left-looking
descending pass at 37$^\circ$ incidence, D141's geometry. The ionosphere, at $\fA$, is a 0.6\,rad/km gradient plus a
Gaussian-correlated field (rms 4\,rad, correlation length 15\,km), about 80\,rad across the scene.

Each band's coherence is a far-field value (0.6, or 0.35) times a damage zone, $q + (1-q)(1-e^{-(d/w)^2})$ with
$q\gamma_0 = 0.12$ and $w=2$\,km north of the trace and $q\gamma_0=0.25$, $w=0.8$\,km south (close to D141's
measured coherence), times the fringe-rate loss $|\mathrm{sinc}|$ of the band's own phase gradient \citep{zebker1992decorrelation}. That loss is
taken over the band's resolution cell (6.2\,m of ground range for the main band, 50\,m for the side band, 5.3\,m
along track) and over the 40\,m pixel.

Each pixel's interferometric sum of $L$ looks (48 for the main band, 6 for the side band) is drawn from the complex
Wishart distribution with that coherence and phase, by the Bartlett decomposition \citep{lee1994intensity}. With
$g_1\sim\Gamma(L)$, $g_2\sim\Gamma(L-1)$ and $a\sim\mathcal{CN}(0,1)$,
\begin{equation}
  \begin{aligned}
  W_{12} &= g_1\gamma e^{i\phi} + \sqrt{g_1}\,\bar a\,\sqrt{1-\gamma^2}, \qquad W_{11} = g_1,\\
  W_{22} &= \gamma^2 g_1 + 2\,\mathrm{Re}\!\left(\gamma e^{-i\phi}\sqrt{g_1}\,\bar a\sqrt{1-\gamma^2}\right)
  + (1-\gamma^2)(|a|^2+g_2).
  \end{aligned}
\end{equation}
The pixels are summed over 240\,m cells exactly as GSLC pixels are, which gives 1728 and 216 looks per cell
(\texttt{split\_spectrum}'s own estimate for NISAR's 40\,+\,5\,MHz mode: 1753 and 217). The main band's unwrapping is
the oracle: the cell mean of the true phase, with the noisy cell phase re-wrapped onto it. All estimators therefore
see the same, correct $\phiA$, and differ only in how they use the side band.
